\documentclass[11pt]{article}
\usepackage[T1]{fontenc}
\usepackage[margin=1in]{geometry}
\usepackage{enumitem}
\usepackage{hyperref}
\setlength{\parindent}{1.5em}
\setlength{\parskip}{6pt}
% =====================================================================
%  EDIT YOUR DETAILS HERE — this ONE file feeds every abstract & deck.
%  (Change a value, recompile; all 36 documents update.)
% =====================================================================
\newcommand{\seminarName}{Aflah Muhammed P}      % <-- your full name (guessed from your email; please verify)
\newcommand{\seminarRoll}{TVE00CS000}            % <-- your university register/roll number
\newcommand{\seminarClassRoll}{00}               % <-- your class roll no (the short one)
\newcommand{\seminarGuide}{Prof.\ <Guide Name>}  % <-- your guide
\newcommand{\seminarDept}{Department of Computer Science and Engineering}
\newcommand{\seminarCollege}{College of Engineering Trivandrum}
\newcommand{\seminarShortCollege}{CET}           % shown in the slide footer, e.g. "Aflah Muhammed P (CET)"
\newcommand{\seminarShortTitle}{B.Tech Seminar}  % middle of the slide footer
\newcommand{\seminarDate}{July 31, 2026}
% ---------------------------------------------------------------------
% Optional: put a logo image named  cet_logo.png  in this folder and it
% will appear on every title slide automatically. If absent, it is skipped.
% =====================================================================

\begin{document}
\begin{center}
{\LARGE\bfseries SWE-smith: Scaling Data for Software-Engineering Agents}\\[1.1em]
{\large\bfseries Seminar Abstract}\\[0.6em]
{\normalsize \seminarDate}
\end{center}
\vspace{0.6em}
\section*{Abstract}
Language-model agents are increasingly able to resolve real-world software issues, but their progress is gated by the availability of training data. Benchmarks such as SWE-bench are built by mining task instances from historical GitHub pull requests, a process that is slow, labor-intensive, and storage-hungry: reproducing an execution environment for every instance can demand tens to hundreds of terabytes. As a result, publicly available training sets have been limited to a few thousand instances, capping the performance of open-weight coding agents and leaving a wide gap behind proprietary systems.

This seminar presents \emph{SWE-smith}, an automated data factory that converts any Python repository into hundreds to thousands of verifiable bug-fixing tasks. SWE-smith builds a single execution environment per repository and then synthetically introduces bugs -- through LM rewriting, LM modification, procedural AST edits, and pull-request mirroring -- that break existing passing tests; the untouched original code together with its passing tests forms a self-labeled, execution-verified fix, while an LM drafts a realistic issue description from the diff and failing test. The pipeline yields roughly 50{,}000 task instances from 128 repositories using only about 295 GB, on the order of 500x less storage than comparable pipelines. Fine-tuning Qwen2.5-Coder-32B on 5{,}016 agent trajectories produces \emph{SWE-agent-LM-32B}, which resolves 40.2\% of SWE-bench Verified in a single attempt without inference-time scaling -- a new state of the art among open models. The presentation walks through the pipeline, the bug-generation strategies, the training recipe, and the empirical results and limitations.
\section*{Key Reference Papers}
\begin{enumerate}
  \item J. Yang, K. Lieret, C. E. Jimenez, A. Wettig, K. Khandpur, Y. Zhang, B. Hui, O. Press, L. Schmidt, and D. Yang, ``SWE-smith: Scaling Data for Software Engineering Agents,'' arXiv:2504.21798, 2025.
  \item C. E. Jimenez, J. Yang, A. Wettig, S. Yao, K. Pei, O. Press, and K. Narasimhan, ``SWE-bench: Can Language Models Resolve Real-World GitHub Issues?,'' in Proc. Int. Conf. Learning Representations (ICLR), 2024.
  \item J. Yang, C. E. Jimenez, A. Wettig, K. Lieret, S. Yao, K. Narasimhan, and O. Press, ``SWE-agent: Agent-Computer Interfaces Enable Automated Software Engineering,'' in Proc. Advances in Neural Information Processing Systems (NeurIPS), 2024.
\end{enumerate}
\vspace{0.8em}
\noindent\textbf{Submitted By}\\
\seminarName\\
\seminarRoll

\vspace{0.6em}
\noindent\textbf{Guide}\\
\seminarGuide
\end{document}
